\section*{\texorpdfstring{\S\CurrentExerciseGroup}{Section \CurrentExerciseGroup}}
\phantomsection
\addcontentsline{toc}{section}{\texorpdfstring{Exercises for \S\CurrentExerciseGroup}{Exercises for Section \CurrentExerciseGroup}}

\begin{enumerate}
\def\labelenumi{\arabic{enumi})}
\tightlist
\item
  Let H be a set, directed for the relation \(\leqslant\) , of integrable functions \(\geqslant 0\) such that \(\sup N_{1}(f) < +\infty\) . Let g be the upper envelope of H; in order that g be integrable and \(f \in H\)
\end{enumerate}

that, in \(L^{1}\) , the class \(\widetilde{g}\) of g be the limit of the filter of sections of the directed set of classes of the functions \(f \in H\) , it is necessary and sufficient that, for every \(\varepsilon > 0\) , there exist a function \(f_{1} \in H\) , a set B, directed for \(\leqslant\) , of lower semi-continuous integrable functions, and a mapping \(f \mapsto f^{*}\) of the set \(H_{1}\) of functions \(f \in H\) that are \(\leqslant f_{1}\) into the set B, such that \(f \leqslant f^{*}\) and \(\mathrm{N}_{1}(f^{*} - f) \leqslant \varepsilon\) for every function \(f \in H_{1}\) (make use of Th. 3 of No.~4).

\begin{enumerate}
\def\labelenumi{\arabic{enumi})}
\setcounter{enumi}{1}
\item
  Let \(\mu\) be Lebesgue measure on \(\mathbf{R}\) and let \(\Omega\) be the topological space obtained by equipping the space \(\mathcal{L}^1\) with the topology of pointwise convergence in \(\mathbf{R}\) . Show that the mapping \(f\mapsto \int fd\mu\) of \(\Omega\) into \(\mathbf{R}\) is not continuous at any point of \(\Omega\) .
\item
  Let I be an interval in \(\mathbf{R}\) , \(\mathbf{f}\) a mapping of \(X \times I\) into a Banach space \(F\) , such that: \(1^{\circ}\) for every \(\alpha \in I\) , the mapping \(t \mapsto \mathbf{f}(t, \alpha)\) of \(X\) into \(F\) is integrable; \(2^{\circ}\) for every \(t \in X\) , the mapping \(\alpha \mapsto \mathbf{f}(t, \alpha)\) has a derivative \(\mathbf{f}_{\alpha}'(t, \alpha)\) in \(I\) ; \(3^{\circ}\) there exists an integrable function \(g \geqslant 0\) such that \(|\mathbf{f}_{\alpha}'(t, \alpha)| \leqslant g(t)\) for all \(t \in X\) and all \(\alpha \in I\) . Under these conditions, show that the function \(\mathbf{u}(\alpha) = \int \mathbf{f}(t, \alpha) d\mu(t)\) is differentiable in \(I\) and that
\end{enumerate}

\[
\mathbf {u} ^ {\prime} (\alpha) = \int \mathbf {f} _ {\alpha} ^ {\prime} (t, \alpha) d \mu (t).
\]

\begin{enumerate}
\def\labelenumi{\arabic{enumi})}
\setcounter{enumi}{3}
\tightlist
\item
  Let \(\mu\) be Lebesgue measure on the interval \(X = [0,1]\) .
\end{enumerate}

\begin{enumerate}
\def\labelenumi{\alph{enumi})}
\item
  Define in X a nowhere dense perfect set A whose measure is an arbitrary number \(\alpha\) such that \(0 \leqslant \alpha < 1\) (use the method of construction of Cantor's triadic set).
\item
  Define in X a sequence \((\mathbf{A}_n)\) of pairwise disjoint nowhere dense integrable sets, such that \(\mu (\mathrm{A}_n) = 2^{-n}\) and such that every interval contiguous to \(\mathbf{B}_n = \bigcup_{k = 1}^{n}\mathbf{A}_k\) contains a subset of \(\mathrm{A}_{n + 1}\) of measure \(>0\) . If \(\mathbf{A} = \bigcup_{n}\mathbf{A}_n\) , show that \(\mathbf{A}\) is a meager set of measure 1, and \(\mathbb{C}\mathbf{A}\) a negligible non-meager set.
\item
  Let \(\mathrm{H} = \bigcup_{n=0}^{\infty} \mathrm{A}_{2n+1}\) ; show that, for every open interval \(\mathrm{I} \subset \mathrm{X}\) , the intersections of I with H and with CH have measure \textgreater{} 0. Let f be a function equal almost everywhere to the characteristic function \(\varphi_{H}\) ; show that there does not exist any sequence \((f_{n})\) of continuous functions on X, convergent at every point of X, whose limit is equal to f (note that f is necessarily discontinuous at every point of X, and make use of Exer. 22 of GT, IX, §5).
\end{enumerate}

\begin{enumerate}
\def\labelenumi{\arabic{enumi})}
\setcounter{enumi}{4}
\tightlist
\item
  Let \(\mu\) be a positive measure on a locally compact space \(X\) . For every numerical function \(f\) (finite or not), of whatever sign, defined on \(X\) , we denote by \(\mu^{*}(f)\) (and call upper integral of \(f\) ) the infimum of the numbers \(\mu(h)\) for the functions \(h \geqslant f\) that are integrable and lower semi-continuous, when such functions exist, and by \(+\infty\) in the contrary case; this definition coincides with that of §1, No.~3 when \(f \geqslant 0\) . We denote by \(\mu_{*}(f)\) , and call lower integral of \(f\) , the number \(-\mu^{*}(-f)\) .
\end{enumerate}

\begin{enumerate}
\def\labelenumi{\alph{enumi})}
\item
  Show that if \(f_{1}\) and \(f_{2}\) are two numerical functions such that \(f_{1}(x) \leqslant f_{2}(x)\) almost everywhere, then \(\mu^{*}(f_{1}) \leqslant \mu^{*}(f_{2})\) .
\item
  Let \(f_{1}\) and \(f_{2}\) be two numerical functions such that \(\mu^{*}(f_{1}) + \mu^{*}(f_{2})\) is defined and \(< +\infty\) ; show that \(f_{1}(x) + f_{2}(x)\) is defined almost everywhere and \(\mu^{*}(f_{1} + f_{2}) \leqslant \mu^{*}(f_{1}) + \mu^{*}(f_{2})\) (reduce to the case that \(f_{1}(x) < +\infty\) and \(f_{2}(x) < +\infty\) at every point, with the help of \(a\) )).
\item
  If \((f_n)\) is an increasing sequence of numerical functions such that \(\mu^{*}(f_{n}) > -\infty\) from some index onward, show that \(\mu^{*}(\sup_{n} f_{n}) = \sup_{n} \mu^{*}(f_{n})\) .
\end{enumerate}

¶ 6) a) Let \(f\) be a numerical function such that \(\mu^{*}(f)\) (Exer. 5) is finite. Show that there exists an integrable function \(f_{1} \geqslant f\) such that \(\mu(f_{1}) = \mu^{*}(f)\) ; if \(f_{2}\) is a second integrable function such that \(f_{2} \geqslant f\) and \(\mu(f_{2}) = \mu^{*}(f)\) , then \(f_{1}\) and \(f_{2}\) are equivalent.

\begin{enumerate}
\def\labelenumi{\alph{enumi})}
\setcounter{enumi}{1}
\item
  For a numerical function \(f\) to be integrable, it is necessary and sufficient that \(\mu^{*}(f)\) and \(\mu_{*}(f)\) be finite and equal.
\item
  Let \(f\) be a numerical function such that \(\mu^{*}(f)\) and \(\mu_{*}(f)\) are both finite; let \(g\) and \(h\) be two integrable functions such that \(g \leqslant f \leqslant h\) and \(\mu(g) = \mu_{*}(f)\) , \(\mu(h) = \mu^{*}(f)\) . Show that
\end{enumerate}

\[
\mu_ {*} (f - g) = \mu_ {*} (h - f) = 0 \quad \text { and } \quad \mu^ {*} (f - g) = \mu^ {*} (h - f) = \mu^ {*} (f) - \mu_ {*} (f).
\]

\begin{enumerate}
\def\labelenumi{\alph{enumi})}
\setcounter{enumi}{3}
\tightlist
\item
  Let \(f_{1}, f_{2}\) be two numerical functions such that the numbers \(\mu^{*}(f_{1}), \mu^{*}(f_{2})\) , \(\mu_{*}(f_{1})\) and \(\mu_{*}(f_{2})\) are all finite; show that
\end{enumerate}

\[
\mu_ {*} (f _ {1} + f _ {2}) \leqslant \mu_ {*} (f _ {1}) + \mu^ {*} (f _ {2}) \leqslant \mu^ {*} (f _ {1} + f _ {2})
\]

(if \(g_2\) is an integrable function such that \(f_2 \leqslant g_2\) and \(\mu^*(f_2) = \mu(g_2)\) , note that, for every integrable function \(h\) such that \(h \leqslant f_1 + f_2\) , one has \(h - g_2 \leqslant f_1\) ). From this, deduce that if \(f_1\) is integrable then

\[
\mu^ {*} (f _ {1} + f _ {2}) = \mu (f _ {1}) + \mu^ {*} (f _ {2}), \quad \mu_ {*} (f _ {1} + f _ {2}) = \mu (f _ {1}) + \mu_ {*} (f _ {2}),
\]

and \(\mu^{*}(f_{1} + f_{2}) = \mu^{*}\big(\sup (f_{1},f_{2})\big) + \mu^{*}\big(\inf (f_{1},f_{2})\big)\) (for the latter relation, make use of Exer. 1 of §1).

\begin{enumerate}
\def\labelenumi{\alph{enumi})}
\setcounter{enumi}{4}
\tightlist
\item
  Let \(f\) be an integrable function. In order that a function \(g\) , such that \(\mu^{*}(g)\) is finite, be integrable, it is necessary and sufficient that \(\mu(f) = \mu^{*}(g) + \mu^{*}(f - g)\) (if \(g_{1}\) is an integrable function such that \(g \leqslant g_{1}\) and \(\mu^{*}(g) = \mu(g_{1})\) , note that \(f - g_{1} \leqslant f - g\) ).
\end{enumerate}

¶ 7) Let \(\mu\) be a positive measure on X. For every subset \(A \subset X\) , the lower integral (Exer. 5) of the characteristic function \(\varphi_A\) is called the inner measure of A and is denoted \(\mu_*(A)\) .

\begin{enumerate}
\def\labelenumi{\alph{enumi})}
\item
  Show that \(\mu_{*}(\mathrm{A})\) is the supremum of the measures of the compact sets contained in A (argue as in Th. 4).
\item
  For every subset A of X of finite outer measure, show that there exist two integrable sets \(A_{1}, A_{2}\) such that \(A_{1} \subset A \subset A_{2}\) and \(\mu_{*}(A) = \mu(A_{1})\) , \(\mu^{*}(A) = \mu(A_{2})\) . For A to be integrable, it is necessary and sufficient that \(\mu^{*}(A)\) and \(\mu_{*}(A)\) be finite and equal. With the same notations, show that
\end{enumerate}

\[
\mu_ {*} (\mathrm{A} \cap \mathbb {C} \mathrm{A} _ {1}) = \mu_ {*} (\mathrm{A} _ {2} \cap \mathbb {C} \mathrm{A}) = 0
\]

and

\[
\mu^ {*} (\mathrm{A} \cap \mathbb {C} \mathrm{A} _ {1}) = \mu^ {*} (\mathrm{A} _ {2} \cap \mathbb {C} \mathrm{A}) = \mu^ {*} (\mathrm{A}) - \mu_ {*} (\mathrm{A}).
\]

\begin{enumerate}
\def\labelenumi{\alph{enumi})}
\setcounter{enumi}{2}
\item
  Let A be an integrable set; show that for every set \(\mathbf{B} \subset \mathbf{A}\) , one has \(\mu(\mathbf{A}) = \mu^{*}(\mathbf{B}) + \mu_{*}(\mathbf{A} \cap \mathbb{C}\mathbf{B})\) .
\item
  Let A and B be two sets of finite outer measure that are disjoint. Show that if \(\mathbf{C} = \mathbf{A} \cup \mathbf{B}\) , then
\end{enumerate}

\[
\mu_ {*} (\mathrm{A}) + \mu_ {*} (\mathrm{B}) \leqslant \mu_ {*} (\mathrm{C}) \leqslant \mu_ {*} (\mathrm{A}) + \mu^ {*} (\mathrm{B}) \leqslant \mu^ {*} (\mathrm{C}) \leqslant \mu^ {*} (\mathrm{A}) + \mu^ {*} (\mathrm{B})
\]

and that

\[
\mu_ {*} (\mathrm{C}) - \mu_ {*} (\mathrm{A}) - \mu_ {*} (\mathrm{B}) \leqslant \mu^ {*} (\mathrm{A}) + \mu^ {*} (\mathrm{B}) - \mu^ {*} (\mathrm{C})
\]

(for the latter inequality, reduce to the case that \(\mu_{*}(\mathrm{A}) = \mu_{*}(\mathrm{B}) = 0\) with the help of \(b\) ); if \(\mathrm{A}_2\) and \(\mathrm{B}_2\) are integrable sets such that \(\mathrm{A} \subset \mathrm{A}_2\) , \(\mathrm{B} \subset \mathrm{B}_2\) , \(\mu^{*}(\mathrm{A}) = \mu(\mathrm{A}_2)\) , \(\mu^{*}(\mathrm{B}) = \mu(\mathrm{B}_2)\) , show that \(\mu_{*}(\mathrm{C}) \leqslant \mu(\mathrm{A}_2 \cap \mathrm{B}_2)\) .

¶ 8) If the torus T = R/Z is canonically identified with the interval {[}0,1 {[} of R, the Lebesgue measure μ on this interval is a measure on T; for every set A ⊂ T and every z ∈ T, μ*(A + z) = μ*(A).

\begin{enumerate}
\def\labelenumi{\alph{enumi})}
\item
  Show that there exists a subgroup \(\mathrm{H}_0\) of \(\mathbf{T}\) such that the sets \(\mathrm{H}_n = r_n + \mathrm{H}_0\) , where \(r_n\) runs over the set of rational numbers contained in {[}0,1{[}, form a partition of \(\mathbf{T}\) (regarding \(\mathbf{R}\) as a vector space over \(\mathbf{Q}\) , note that in \(\mathbf{R}\) the subspace \(\mathbf{Q}\) admits a supplement).
\item
  Let H be a set that is the union of any finite number of sets \(H_{n}\) ; show that H is not integrable and that \(\mu_{*}(H)=0\) (note that every subgroup of the additive group Q/Z generated by a finite number of elements has infinite index in Q/Z; from this, deduce that there exists a partition of T into a countable infinity of sets \(P_{n}\) , each of which contains a set of the form \(z+H\) ).
\item
  Deduce from \(b\) ) an example of a decreasing sequence \((\mathrm{A}_n)\) of subsets of \(\mathbf{T}\) , whose intersection is empty, and is such that \(\mu^{*}(\mathrm{A}_{n}) = 1\) for all \(n\) .
\item
  Let A be an integrable set such that \(H_{0} \subset A\) and \(\mu(A) = \mu^{*}(H_{0}) > 0\) ; show that, for every integrable set \(B \subset A\) , the sets \(B_{1} = B \cap H_{0}\) and \(B_{2} = B \cap C H_{0}\) form a partition of B such that \(\mu^{*}(B_{1}) = \mu^{*}(B_{2}) = \mu(B)\) and \(\mu_{*}(B_{1}) = \mu_{*}(B_{2}) = 0\) (make use of Exer. 7 b)).
\end{enumerate}

¶ 9) a) Let \(\mu\) be a positive measure on a locally compact space \(X\) . In the Banach space \(\widetilde{\mathcal{F}}^1\) of equivalence classes of functions in \(\mathcal{F}^1\) , let \(G\) be a closed linear subspace containing the subspace \(L^{1}\) ; let G be the closed linear subspace of \(F^{1}\) formed by the functions whose class belongs to G. Show that one can extend to G the linear form \(\mu(f)\) in such a way that the inequality \(|\mu(f)| \leqslant N_{1}(f)\) is again verified (make use of the Hahn--Banach theorem). Show that for the integral so extended, Th. 3 of §3 is again valid.

\begin{enumerate}
\def\labelenumi{\alph{enumi})}
\setcounter{enumi}{1}
\tightlist
\item
  Assume that G is the direct sum of \(\mathbf{L}^1\) and a subspace H of finite dimension. Show that, in \(\mathcal{G}\) , Lebesgue's theorem (Th. 2) is again valid. (Let \((f_n)\) be a sequence of functions in \(\mathcal{G}\) , tending almost everywhere to \(f\) and such that \(|f_n| \leqslant g\) , where \(g \geqslant 0\) and \(\mathrm{N}_1(g) < +\infty\) . Let \((\widetilde{u}_k)_{1 \leqslant k \leqslant m}\) be a basis of H, and let \(f_n = \sum_{k=1}^{m} \alpha_{nk} u_k + h_n\) , where \(h_n \in \mathcal{L}^1\) . Using the fact that \(G\) is the topological direct sum of \(L^1\) and \(H\) , show that the \(\alpha_{nk}\) are uniformly bounded; on passing, if necessary, to a subsequence of \((f_n)\) , reduce to the case that each of the sequences \((\alpha_{nk})_{n \geqslant 1}\) has a limit; from this, deduce that \(h_n(x)\) then tends to a limit almost everywhere, and apply Lebesgue's theorem to the sequence \((h_n)\) ; note finally that two subsequences of \((\widetilde{f}_n)\) cannot tend to different limits in \(G\) .
\end{enumerate}

¶ 10) a) Let E be a Riesz space, μ a positive linear form on E such that the relation μ(\textbar x\textbar) = 0 implies x = 0; μ(\textbar x\textbar) is then a norm on E, and let us assume that E, equipped with this norm, is complete. It then follows that E is fully lattice-ordered (Ch. II, §2, Exer. 8 e)). Consequently (Ch. II, §1, Exers. 12 and 13) there exists a locally compact space X, the sum of a family (Kα) of compact Stone spaces, such that E is isomorphic to a space of continuous numerical functions (finite or not) on X that contains the space K(X). Identifying E with this space, the restriction of μ to K(X) is then a positive measure on X. Show that E is canonically isomorphic to the space L¹(μ); more precisely, for every μ-integrable function g defined on X, there exists one and only one function f ∈ E that is equivalent to g for the measure μ (note that every element ≥ 0 of E is the supremum of an increasing sequence of elements of K(X), and that K(X) is dense in L¹(μ)).

\begin{enumerate}
\def\labelenumi{\alph{enumi})}
\setcounter{enumi}{1}
\tightlist
\item
  Deduce from a) that, for every compact space K, there exist a locally compact space S, the topological sum of a family of compact Stone spaces, and a positive measure \(\nu\) on S, with support equal to S, such that the fully lattice-ordered space \(\mathcal{M}(\mathrm{K})\) of measures on K, equipped with the norm \(\|\mu\|\) , is isomorphic to \(\mathcal{L}^{1}(\nu)\) (consider the linear form \(\mu\mapsto\mu(\mathrm{K})\) on \(\mathcal{M}(\mathrm{K})\) ).
\end{enumerate}

\begin{enumerate}
\def\labelenumi{\arabic{enumi})}
\setcounter{enumi}{10}
\tightlist
\item
  \begin{enumerate}
  \def\labelenumii{\alph{enumii})}
  \tightlist
  \item
    Let \(\Gamma\) be any set of subsets of a set \(A\) . Let \(\Psi\) the set of subsets of \(A\) of the form
  \end{enumerate}
\end{enumerate}

\[
\mathrm{X} _ {1} \cap \mathrm{X} _ {2} \cap \dots \cap \mathrm{X} _ {m} \cap \mathbb {C} \mathrm{X} _ {m + 1} \cap \dots \cap \mathbb {C} \mathrm{X} _ {m + p},
\]

where the \(X_{i}\) are sets in \(\Gamma\) , m is any integer \(\geqslant 1\) , and p is any integer \(\geqslant 0\) . Show that the smallest clan \(\Phi\) containing \(\Gamma\) is the set of finite unions of sets in \(\Psi\) .

\begin{enumerate}
\def\labelenumi{\alph{enumi})}
\setcounter{enumi}{1}
\item
  Let \(\Delta\) be the set of finite intersections of sets in \(\Gamma\) . Show that, for every vector space F, the set of linear combinations (with coefficients in F) of characteristic functions of sets of the clan generated by \(\Gamma\) is identical to the set of linear combinations of characteristic functions of sets in \(\Delta\) .
\item
  Let X be a topological space, \(\Gamma\) the set of compact subsets of X. Show that the clan generated by \(\Gamma\) is identical to the set of finite unions of subsets of X of the form \(A \cap CB\) , where A and B are compact sets.
\end{enumerate}

\begin{enumerate}
\def\labelenumi{\arabic{enumi})}
\setcounter{enumi}{11}
\tightlist
\item
  Let \(X\) be a Hausdorff topological space. For every pair \((K, U)\) formed by a compact set \(K\) and an open set \(U\) in \(X\) , let \(I(K, U)\) be the set of subsets \(M \subset X\) such that \(K \subset M \subset U\) , and let \(\mathcal{T}\) be the topology on \(\mathfrak{P}(X)\) generated by the set of subsets \(I(K, U)\) of \(\mathfrak{P}(X)\) .
\end{enumerate}

\begin{enumerate}
\def\labelenumi{\alph{enumi})}
\item
  Show that each of the sets I(K, U) is both open and closed in \(\mathfrak{P}(X)\) ; from this, deduce that \(\mathfrak{P}(X)\) , equipped with the topology \(\mathcal{T}\) , is a completely regular and totally disconnected space.
\item
  For the mapping \(M \mapsto CM\) of \(\mathfrak{P}(X)\) onto itself to be continuous (for the topology T), it is necessary and sufficient that X be compact.
\item
  Take X to be the interval {[}0,1{]} of R. Show that the mappings \((M,N)\mapsto M\cup N\) and \((M,N)\mapsto M\cap N\) of \(\mathfrak{P}(X)\times\mathfrak{P}(X)\) into \(\mathfrak{P}(X)\) are not continuous for the topology T.
\item
  Assume X to be locally compact. Show that the topology induced by T on the set of compact subsets of X is finer than the topology deduced from a uniform structure on X by the procedure of Exer. 5 of GT, II, §1; these two topologies cannot be identical unless X is a discrete space.
\end{enumerate}

¶ 13) a) Let X be a locally compact space, and let \(\Gamma\) be a base for the topology of X consisting of relatively compact sets. Let \(\mathcal{U}\) be a uniform structure compatible with the topology of X, and let \(\mathfrak{S}\) be a fundamental system of entourages for this structure. Let \(M \mapsto \lambda(M)\) be a finite numerical function \(\geqslant 0\) defined on \(\Gamma\) . For every compact set \(K \subset X\) and every entourage \(V \in \mathfrak{S}\) , let \(\alpha_V(K)\) be the infimum of the numbers \(\sum_{i} \lambda(U_i)\)

for all finite coverings \((\mathrm{U}_{i})\) of K formed by sets in \(\Gamma\) that are small of order V; suppose that as V runs over S, the supremum \(\alpha(\mathrm{K})\) of the numbers \(\alpha_{\mathrm{V}}(\mathrm{K})\) is finite. Show that there exists a measure \(\mu\) on X (and only one) such that \(\mu(\mathrm{K}) = \alpha(\mathrm{K})\) for every compact set K (make use of Th. 5).

\begin{enumerate}
\def\labelenumi{\alph{enumi})}
\setcounter{enumi}{1}
\tightlist
\item
  Let \(\Psi\) be the set of Borel subsets of \(X\) , \(\alpha\) a mapping of \(\Psi\) into \([0, +\infty]\) satisfying the following conditions:
\end{enumerate}

\begin{enumerate}
\def\labelenumi{(\roman{enumi})}
\item
  If \(\mathrm{B}_1, \mathrm{B}_2\) are two disjoint Borel subsets of \(X\) , then \(\alpha(\mathrm{B}_1 \cup \mathrm{B}_2) = \alpha(\mathrm{B}_1) + \alpha(\mathrm{B}_2)\) .\\
\item
  If \(B\) is a compact subset of \(X\) then \(\alpha(B) < +\infty\) .
\item
  If \(\mathbf{B}\) is a Borel subset of \(\mathbf{X}\) then \(\alpha(\mathbf{B}) = \inf \alpha(\mathbf{U})\) , where \(\mathbf{U}\) runs over the set of open subsets of \(\mathbf{X}\) containing \(\mathbf{B}\) .
\end{enumerate}

Then, there exists one and only one positive measure \(\mu\) on X such that \(\alpha(B) = \mu^{*}(B)\) for every Borel subset B of X that can be covered by a sequence of compact sets. c) For every \(B \in \Phi\) , set \(\alpha(B) = 0\) if B can be covered by a sequence of compact sets, and \(\alpha(B) = +\infty\) in the contrary case. Then, the conditions (i), (ii), (iii) of \(b\) are satisfied. One has \(\mu = 0\) , therefore \(\alpha(B) \neq \mu^{*}(B)\) if B cannot be covered by a sequence of compact sets.

\begin{enumerate}
\def\labelenumi{\arabic{enumi})}
\setcounter{enumi}{13}
\item
  Let \((\mathbf{A}_n)\) be a sequence of integrable sets such that \(\sum_{n} \mu(\mathbf{A}_n) < +\infty\) . For every integer \(k\) , let \(\mathbf{G}_k\) be the set of \(x \in X\) such that \(x \in \mathbf{A}_n\) for at least \(k\) values of \(n\) ; show that \(\mathbf{G}_k\) is integrable and that \(k \cdot \mu(\mathbf{G}_k) \leqslant \sum_{n} \mu(\mathbf{A}_n)\) .
\item
  Let \(\mu\) be a bounded positive measure on a locally compact space \(X\) , and let \((A_n)\) be a sequence of integrable sets in \(X\) such that \(\inf_{n} \mu(A_n) = m > 0\) . Show that the set \(B\) of points of \(X\) that belong to infinitely many of the sets \(A_n\) is integrable, and that \(\mu(B) \geqslant m\) .
\end{enumerate}

¶ 16) Let X be a completely regular space, and let \(\mathcal{C}(X)\) (resp. \(\mathcal{C}^b(X)\) ) be the Riesz space of continuous (resp. continuous and bounded) numerical functions on X.

\begin{enumerate}
\def\labelenumi{\alph{enumi})}
\item
  Show that if a linear form \(\lambda\) on \(\mathcal{C}(\mathrm{X})\) is continuous for the topology of compact convergence, it is relatively bounded.
\item
  Let \(\lambda\) be a positive linear form on \(\mathcal{C}(\mathrm{X})\) . Show that if \(\lambda\) is zero on \(\mathcal{C}^b (\mathrm{X})\) , then it is zero on \(\mathcal{C}(\mathrm{X})\) (let \(\varphi\) be the canonical mapping of \(\mathcal{C}(\mathrm{X})\) onto the quotient Riesz space \(\mathcal{C}(\mathrm{X}) / \mathcal{C}^b (\mathrm{X})\) (Ch. II, §1, Exer. 4); show that if \(h\) is a continuous function \(\geqslant 0\) and is not bounded on \(\mathrm{X}\) , then \(n\varphi (h)\leqslant \varphi (h^{2})\) for every integer \(n > 0\) ).
\item
  Let \(\beta X\) be the Stone--Čech compactification of \(X\) , the compact space obtained by equipping \(X\) with the coarsest uniform structure for which the functions in \(\mathcal{C}^b(X)\) are uniformly continuous and by completing the uniform space so obtained; every function \(f \in \mathcal{C}(X)\) may then be extended by continuity to a continuous function \(\widetilde{f}\) on \(\beta X\) , with possibly infinite values (consider \(f/(1 + |f|)\) ). If \(\lambda\) is a positive linear form on \(\mathcal{C}(X)\) , then the restriction of \(\lambda\) to \(\mathcal{C}^b (\mathrm{X})\) is of the form \(f\mapsto \mu (\widetilde{f})\) , where \(\mu\) is a positive measure on \(\beta X\) ; show that for every function \(f\in \mathcal{C}(\mathrm{X})\) , \(\widetilde{f}\) is integrable for \(\mu\) and \(\lambda (f) = \mu (\widetilde{f})\) (make use of \(b\) ), noting that every function \(\geqslant 0\) belonging to \(\mathcal{C}(\mathrm{X})\) is the upper envelope of a sequence of functions in \(\mathcal{C}^b (\mathrm{X})\) ).
\item
  Show that every point \(x_{0}\) of the support of \(\mu\) that does not belong to X has the following property: for every decreasing sequence \((V_{n})\) of neighborhoods of \(x_{0}\) , the intersection of the \(V_{n}\) contains at least one point of X. Converse.
\item
  Deduce from \(d\) ) that if \(X\) is locally compact and countable at infinity, then the support of \(\mu\) is contained in \(X\) (compare with \(h\) )).
\item
  If X is discrete, show that the support of \(\mu\) is finite (in the contrary case, form a function \(f \geqslant 0\) defined on X such that \(\widetilde{f}\) is not \(\mu\) -integrable).
\item
  Show that if a linear form \(\lambda\) on \(\mathcal{C}(\mathbf{X})\) is positive and continuous for the topology of compact convergence, then the support of \(\mu\) is contained in X (cf.~Ch. III, §2, No.~3, Prop. 11).
\item
  Let \(X_{0}\) be the compact space obtained by adjoining a point at infinity \(\omega\) to the locally compact space X defined in GT, I, §9, Exer. 11, let Y be the subspace of \(\overline{R}\) formed by the integers \(\geqslant 0\) and \(+\infty\) , and let Z be the locally compact space complementary to the point \((\omega, +\infty)\) in the product space \(X_{0} \times Y\) . Show that every continuous numerical function on Z is bounded and that \(f(z)\) tends to a limit as z tends to the point at infinity \((\omega, +\infty)\) of Z. From this, deduce that if \(\mu\) is the measure on Z defined by the mass \(1/2^{n}\) placed at the point \((\omega, n)\) for every \(n \geqslant 0\) , then every continuous function on Z is \(\mu\) -integrable, but the support of \(\mu\) is not compact.
\end{enumerate}

¶ 17) Let X be the locally compact space defined in GT, I, §9, Exer. 11.

\begin{enumerate}
\def\labelenumi{\alph{enumi})}
\tightlist
\item
  Let H be a compact subset of the locally convex space \(\mathcal{C}(\mathrm{X};\mathbf{C})\) equipped with the topology of compact convergence; show that the functions \(f \in H\) are uniformly bounded on X and that there exists \(c \in X\) such that, for every \(x \geqslant c\) , all of the functions of H are constant. (Argue by contradiction; if the functions of H were not uniformly bounded on X, there would be an increasing sequence \((x_{n})\) of points of X such that \(\sup(|\mathrm{H}(x_{n})|) \geqslant n\) , where \(|\mathrm{H}(x_{n})|\) denotes the set of \(|f(x_{n})|\) for \(f \in H\) ; observe that the sequence \((x_{n})\) is convergent in X. Similarly, for every \(x \in X\) , let \(\delta(x)\) be the supremum of the oscillations of all the functions \(f \in H\) in the interval \([x, \rightarrow[; \delta(x)\) is finite and decreasing, therefore there exists \(d \in X\) such that \(\delta(x)\) is constant for \(x \geqslant d\) . To see that this constant \(\beta\) is zero, argue by contradiction as before, on forming two increasing sequences \((s_{n}), (t_{n})\) of points of X such that \(s_{n} \leqslant t_{n} \leqslant s_{n+1} \leqslant t_{n+1}\) and a sequence \((f_{n})\) of functions in H such that \(|f_{n}(s_{n}) - f_{n}(t_{n})| \geqslant \beta/2\) .
\end{enumerate}

In particular, \(\mathcal{C}(\mathrm{X};\mathbf{C})\) is the direct sum of \(\mathcal{K}(\mathrm{X};\mathbf{C})\) and \(C\cdot1\) .

\begin{enumerate}
\def\labelenumi{\alph{enumi})}
\setcounter{enumi}{1}
\item
  Show that on X, every measure has compact support (for every \(x \in X\) let \(f_x\) be the characteristic function of the interval \(\leftarrow, x\) {]}, which is continuous and has compact support; consider on X the increasing function \(|\mu|(f_x)\) ).
\item
  Deduce from a) and b) that in \(\mathcal{M}(\mathrm{X};\mathbf{C}) = \mathcal{M}^{1}(\mathrm{X};\mathbf{C}) = \mathcal{C}'(\mathrm{X};\mathbf{C})\) , the bounded sets for the topology \(\mathcal{T}\) of uniform convergence on the compact subsets of \(\mathcal{C}(\mathrm{X};\mathbf{C})\) are the bounded sets for the ultrastrong topology (Ch. III, §1, Exer. 15). Show that \(\mathcal{C}'(\mathrm{X};\mathbf{C})\) is not quasi-complete for the topology \(\mathcal{T}\) (consider the measures \(\varepsilon_{x}\) for \(x\in \mathrm{X}\) ); from this, deduce that for the topology of compact convergence, \(\mathcal{C}(\mathrm{X};\mathbf{C})\) is neither barreled nor bornological (cf.~TVS, III, §4, No.~2, Cor. 4 of Th. 1 and Exer. 18 below).
\item
  Show that on \(\mathcal{K}(\mathrm{X};\mathbf{C})\) the topology \(\mathcal{T}_0\) (the direct limit of the topologies of the Banach spaces \(\mathcal{K}(\mathrm{X},\mathrm{K};\mathbf{C})\) ) is identical to the topology of uniform convergence, and that, equipped with this topology, \(\mathcal{K}(\mathrm{X};\mathbf{C})\) is complete. (Let V be a neighborhood of 0 for \(\mathcal{T}_0\) ; for every \(x\in X\) , let \(r_x\) be the largest number \(>0\) such that V contains all of the continuous functions \(f\) with support contained in \([ \leftarrow ,x]\) and such that \(\| f\| \leqslant r_x\) . Show that the infimum of \(r_x\) in X is \(>0\) , by proving that there cannot exist an increasing sequence \((x_n)\) of points of X such that \(\lim_{n\to \infty}r_{x_n} = 0\) .)
\end{enumerate}

\begin{enumerate}
\def\labelenumi{\arabic{enumi})}
\setcounter{enumi}{17}
\tightlist
\item
  Let \(X\) be a paracompact locally compact space.
\end{enumerate}

\begin{enumerate}
\def\labelenumi{\alph{enumi})}
\item
  Show that the space \(\mathcal{C}(\mathrm{X};\mathbf{C})\) is isomorphic to a product of Fréchet spaces, hence is barreled.
\item
  For a subset H of \(\mathcal{C}^{\prime}(X; \mathbf{C})\) to be bounded for the topology of compact convergence, it is necessary and sufficient that there exist a compact subset K of X such that \(\operatorname{Supp}(\mu) \subset K\) for all \(\mu \in H\) and such that H is bounded for the ultrastrong topology (compare with Exer. 17 c)). The topology induced on H by the topology of compact convergence is then identical to the vague topology.
\item
  Show that on the set \(\mathcal{M}_{+}(\mathrm{X})\cap \mathcal{C}'(\mathrm{X};\mathbf{C})\) the topology induced by the topology of compact convergence on \(\mathcal{C}'(\mathrm{X};\mathbf{C})\) is identical to the vague topology. Is the same true when \(X\) is not paracompact (Exer. 17)?
\end{enumerate}

\begin{enumerate}
\def\labelenumi{\arabic{enumi})}
\setcounter{enumi}{18}
\tightlist
\item
  Equip \(\mathcal{C}(\mathrm{X};\mathbf{C})\) with the topology of compact convergence, and \(\mathcal{C}'(\mathrm{X};\mathbf{C})\) with the topology of uniform convergence on the compact subsets of \(\mathcal{C}(\mathrm{X};\mathbf{C})\) . Let \(\mathfrak{S}\) be the set of compact subsets of \(\mathcal{C}(\mathrm{X};\mathbf{C})\) ; show that the bilinear mapping \((g,\mu)\mapsto g\cdot \mu\) of \(\mathcal{C}(\mathrm{X};\mathbf{C})\times \mathcal{C}'(\mathrm{X};\mathbf{C})\) into \(\mathcal{C}'(\mathrm{X};\mathbf{C})\) is \(\mathfrak{S}\) -hypocontinuous. Let \(\mathcal{T}\) be the set of bounded subsets of \(\mathcal{C}'(\mathrm{X};\mathbf{C})\) ; show that if X is paracompact, the preceding bilinear mapping is also \(\mathcal{T}\) -hypocontinuous; is the latter property again true when X is no longer assumed to be paracompact (Exer. 17)?
\end{enumerate}

¶ 20) Let X,Y be two paracompact locally compact spaces. Equip the spaces \(\mathcal{C}'(X; \mathbf{C})\) and \(\mathcal{C}'(Y; \mathbf{C})\) with the topology of compact convergence.

\begin{enumerate}
\def\labelenumi{\alph{enumi})}
\item
  Show that when \(\mathcal{C}^{\prime}(X\times Y;\mathbf{C})\) is equipped with the topology of compact convergence, the mapping \((\mu,\nu)\mapsto\mu\otimes\nu\) is \((\mathfrak{S},\mathfrak{T})\) -hypocontinuous, where S denotes the set of bounded subsets of \(\mathcal{C}^{\prime}(X;\mathbf{C})\) and T the set of bounded subsets of \(\mathcal{C}^{\prime}(Y;\mathbf{C})\) (make use of Exer. 18 b)).
\item
  Show that when X and Y are countable at infinity and \(\mathcal{C}^{\prime}(X\times Y;C)\) is equipped with the topology of compact convergence, the mapping \((\mu,\nu)\mapsto\mu\otimes\nu\) is continuous. (Inspired by the proof of Exer. 3 of Ch. III, §4 and introducing in X (resp. Y) a continuous partition of unity subordinate to a locally finite open covering, argue as in TVS, II, §4, Exer. 9 a).
\item
  If X is discrete, the space \(\mathcal{C}^{\prime}(X;\mathbf{C})\) may be identified with the direct sum space \(\mathbf{C}^{(X)}\) , equipped with the finest locally convex topology. From this, deduce an example where X and Y are discrete, with X uncountable, and where the mapping \((\mu,\nu)\mapsto\mu\otimes\nu\) is not continuous when \(\mathcal{C}^{\prime}(X\times Y;\mathbf{C})\) is equipped with the vague topology (cf.~TVS, II, §4, Exer. 9 b)).
\end{enumerate}

\begin{enumerate}
\def\labelenumi{\arabic{enumi})}
\setcounter{enumi}{20}
\item
  Let \(\mathbf{X}\) be a locally compact space, \(f\) a numerical function \(\geqslant 0\) defined on \(\mathbf{X}\) . In order that the mapping \(\mu \mapsto \mu^{*}(f)\) of \(\mathcal{M}_{+}(\mathbf{X}) \cap \mathcal{C}'(\mathbf{X};\mathbf{C})\) into \(\overline{\mathbf{R}}\) be continuous for the topology of compact convergence, it is necessary and sufficient that \(f\) be continuous on \(\mathbf{X}\) .
\item
  Let X, Y be two locally compact spaces, \(\mu\) a bounded measure on X, \(\nu\) a bounded measure on Y, so that \(\mu \otimes \nu\) is a bounded measure on \(X \times Y\) . Show that for every function \(f \in \mathcal{C}^{0}(X \times Y; \mathbf{C})\) , the function \(x \mapsto \int f(x, y) d\nu(y)\) belongs to \(\mathcal{C}^{0}(X; \mathbf{C})\) and
\end{enumerate}

\[
\int f (x, y) d \mu (x) d \nu (y) = \int d \mu (x) \int f (x, y) d \nu (y).
\]

¶ 23) On the space \(\mathbf{R}^n\) , let \(\mu\) be Lebesgue measure, \(|\mathbf{x}|\) a norm such that the unit ball for this norm has measure equal to 1, and \((\mathbf{x}_k)_{k \geqslant 1}\) an infinite sequence of distinct points of a bounded integrable set B such that \(\mu(B) = 1\) . For every integer \(m\) , denote by \(d_m\) the infimum of the numbers \(|\mathbf{x}_i - \mathbf{x}_j|\) for \(1 \leqslant i < j \leqslant m\) . Show that

\(\liminf_{m \to \infty} md_m^n \leqslant \alpha_n^{-1}\) , where

\[
\alpha_ {n} = 1 + n \int_ {0} ^ {1} \frac {(1 - t) ^ {n}}{1 + t} d t.
\]

(Argue by contradiction, assuming that for some \(\varepsilon > 0\) there exists an \(m_0\) such that \(md_m^n > h_n\) for \(m \geqslant m_0\) , with \(h_n = \alpha_n^{-1} + \varepsilon\) . For \(1 \leqslant i < m\) , let \(B_i\) be the ball with center \(\mathbf{x}_i\) and radius \(\frac{1}{2} h_n m^{-1/n}\) , and for \(m \leqslant i \leqslant 2^n m\) let \(B_i\) be the ball with center \(\mathbf{x}_i\) and radius \(\frac{1}{2} h_n (2i^{-1/n} - m^{-1/n})\) . Show that the \(2^n m\) balls \(B_i\) are pairwise disjoint, and evaluate the measure of their union, making use of the Euler--Maclaurin formula.)
% semantic-fragments: legacy-input-seam
\relax{}
\setcounter{exercisegroup}{5}
\edef\CurrentExerciseGroup{\arabic{exercisegroup}}
